\section{Elements of Experimental Design and Analysis of Variance (ANOVA)}

\subsection{Fundamental Concepts of Experimental Design}

In applied statistics and scientific research, observational studies often fail to establish causal relationships due to the presence of confounding variables and uncontrolled environmental noise. \textbf{Design of Experiments (DOE)} provides a rigorous mathematical and structural methodology to plan, execute, and analyze empirical tests where intentional changes are made to input variables (factors) of a system or process to observe, measure, and identify the corresponding changes in output variables (responses).

The primary objectives of experimental design include:
\begin{enumerate}
	\item Maximizing the information obtained regarding the relationship between independent and dependent variables while minimizing experimental cost, time, and sample size.
	\item Ensuring that the estimated experimental error ($\sigma^2$) is unbiased, consistent, and statistically independent of the treatments applied.
	\item Providing exact probabilistic foundations to test hypotheses regarding treatment differences with maximum statistical power.
\end{enumerate}

To guarantee the validity of any statistical inference derived from an experiment, the experimental protocol must strictly satisfy the \textbf{Three Fundamental Principles of Experimental Design}, originally formalized by Sir Ronald A. Fisher:

\subsubsection{1. Replication}
\textbf{Replication} refers to the independent repetition of the basic experiment under identical operational conditions across distinct experimental units ($n \geq 2$). Replication serves two vital mathematical functions:
\begin{itemize}
	\item It permits the calculation of an \textbf{unbiased estimate of experimental error variance ($\sigma^2$)}, which represents the baseline biological, physical, or measurement variation between experimental units treated alike. Without replication ($n=1$), the error sum of squares ($\text{SCE}$) has zero degrees of freedom ($\nu_E = 0$), rendering exact parametric significance testing via the $F$-distribution impossible.
	\item It increases the precision of treatment mean estimators ($\bar{Y}_{i.}$) by reducing the standard error of the sample mean ($\sigma_{\bar{Y}} = \frac{\sigma}{\sqrt{n}}$), directly enhancing the statistical power of the test ($1 - \beta$).
\end{itemize}

\subsubsection{2. Randomization}
\textbf{Randomization} is the random assignment of treatments to experimental units, ensuring that every unit has an equal and independent probability of receiving any given treatment. Mathematically, randomization fulfills several crucial requirements:
\begin{itemize}
	\item It validates the foundational structural assumption that the random error terms ($\varepsilon_{ij}$) are independent and identically distributed ($\text{i.i.d.}$) random variables with zero mean and constant variance.
	\item It transforms unknown, systematic bias or confounding gradients (e.g., spatial variation across a laboratory, temporal drift in instrument calibration) into random, unbiased experimental noise.
	\item It guarantees the validity of parametric tests ($t$-test, $F$-test) by ensuring that the observed differences in sample means are attributable purely to treatment effects or pure chance, rather than pre-existing systemic differences among experimental units.
\end{itemize}

\subsubsection{3. Blocking (Local Control)}
\textbf{Blocking} is the systematic grouping of experimental units into relatively homogeneous sets (called \textbf{blocks}) before assigning treatments. The fundamental philosophy of blocking is to isolate known, predictable sources of nuisance variability across experimental units so that they do not inflate the experimental error.

By including the block effect ($\beta_j$) explicitly in the additive linear model, the variability among blocks ($\text{SCB}$) is orthogonally extracted and removed from the total residual error sum of squares ($\text{SCE}$). This reduction in residual error ($\text{CME}$) significantly tightens confidence intervals around treatment means and dramatically increases the sensitivity and statistical power of the $F$-test to detect subtle treatment effects. As a general rule of experimental design: *"Block what you can systematically identify and control; randomize what you cannot."*

\subsection{One-Way Analysis of Variance (One-Way ANOVA)}

The \textbf{One-Way Analysis of Variance (Completely Randomized Design, CRD)} is an inferential methodology designed to simultaneously compare the population means of $k \geq 2$ independent groups or levels of a single categorical independent variable (called a \textbf{factor}).

While the two-sample Student's $t$-test is adequate for comparing $k=2$ population means, attempting to compare $k > 2$ means using $\binom{k}{2}$ multiple pairwise $t$-tests inevitably leads to a massive inflation of the \textbf{Family-Wise Error Rate ($\text{FWER}$)}. If $\binom{k}{2} = m$ independent pairwise comparisons are conducted, each at a nominal significance level $\alpha$, the true probability of committing at least one Type I error across the family of tests ($\alpha_{\text{global}}$) grows exponentially according to Boole's inequality and the independence bound:
\begin{align}
	\alpha_{\text{global}} = 1 - (1 - \alpha)^m = 1 - (1 - 0.05)^{10} = 1 - (0.95)^{10} \approx 0.4013 \quad (40.13\%).
\end{align}
ANOVA circumvents this inflation by executing a single global (\textbf{omnibus}) test that partitions the total observed variance into orthogonal components, maintaining the exact experimental Type I error rate at $\alpha$.

\subsubsection{Structural Additive Linear Model}

Let $Y_{ij}$ represent the continuous response observation from the $j$-th experimental unit assigned to the $i$-th treatment level ($i = 1, 2, \dots, k$; $j = 1, 2, \dots, n_i$). The total number of observations across the experiment is $N = \sum_{i=1}^k n_i$. When all group sample sizes are equal ($n_1 = n_2 = \dots = n_k = n$), the design is termed \textbf{balanced}, and $N = kn$.

The observation $Y_{ij}$ is modeled by the parameterization of the effects model:
\begin{align}
	Y_{ij} = \mu + \tau_i + \varepsilon_{ij},
\end{align}
where:
\begin{itemize}
	\item $\mu$ is the overall population grand mean (the theoretical average response across all experimental units and treatments).
	\item $\tau_i = \mu_i - \mu$ is the \textbf{treatment effect} associated with the $i$-th level of the factor, subject to the structural identification constraint $\sum_{i=1}^k n_i \tau_i = 0$ (or $\sum_{i=1}^k \tau_i = 0$ for balanced designs).
	\item $\varepsilon_{ij}$ is the \textbf{random experimental error} associated with observation $(i,j)$, assumed to be independent and identically distributed normal variables:
	\begin{align}
		\varepsilon_{ij} \sim N(0, \sigma^2).
	\end{align}
\end{itemize}

\subsubsection{Statistical Assumptions of ANOVA}

For the parametric $F$-statistic to exactly follow the theoretical Snedecor $F$-distribution, the observational data must strictly satisfy three distributional assumptions:
\begin{enumerate}
	\item \textbf{Normality:} The residuals $\varepsilon_{ij} = Y_{ij} - \bar{Y}_{i.}$ are normally distributed within each treatment group across the population ($Y_{ij} \sim N(\mu_i, \sigma^2)$). This assumption is formally evaluated using residual quantile-quantile (Q-Q) plots, the Shapiro-Wilk test, or the Kolmogorov-Smirnov test. By virtue of the Central Limit Theorem, ANOVA is robust against moderate departures from normality when sample sizes are large and balanced.
	\item \textbf{Homoscedasticity (Homogeneity of Variance):} All $k$ treatment populations possess identical intrinsic error variance ($\sigma_1^2 = \sigma_2^2 = \dots = \sigma_k^2 = \sigma^2$). This is the most critical assumption; severe heteroscedasticity combined with unbalanced sample sizes heavily distorts Type I error rates. Homoscedasticity is formally tested via \textbf{Levene's test}, \textbf{Brown-Forsythe test} (robust against non-normality), or \textbf{Bartlett's test} (highly sensitive to normality violations). If violated, \textbf{Welch's ANOVA} or variance-stabilizing transformations (such as logarithmic $\ln Y$ or square root $\sqrt{Y}$) must be employed.
	\item \textbf{Independence of Errors:} The random error terms $\varepsilon_{ij}$ are mutually independent ($\text{Cov}(\varepsilon_{ij}, \varepsilon_{i'j'}) = 0$ for any $(i,j) \neq (i',j')$). This assumption is guaranteed methodologically through proper \textbf{randomization} during experimental execution. Violation of independence (e.g., temporal autocorrelation or spatial clustering) causes severe underestimation of residual error variance ($\text{CME}$), producing spurious false-positive rejections.
\end{enumerate}

\subsubsection{Partitioning of the Total Sum of Squares}

The fundamental algebraic identity of ANOVA decomposes the total deviation of any observation from the grand mean into two distinct, additive components: the deviation of the group mean from the grand mean (between-group effect), and the deviation of the observation from its respective group mean (within-group error):
\begin{align}
	Y_{ij} - \bar{Y}_{..} = (\bar{Y}_{i.} - \bar{Y}_{..}) + (Y_{ij} - \bar{Y}_{i.}),
\end{align}
where the sample group mean is $\bar{Y}_{i.} = \frac{1}{n_i} \sum_{j=1}^{n_i} Y_{ij}$, and the sample grand mean across all observations is $\bar{Y}_{..} = \frac{1}{N} \sum_{i=1}^k \sum_{j=1}^{n_i} Y_{ij}$.

By squaring both sides of this equation and summing over all $i = 1, \dots, k$ and $j = 1, \dots, n_i$, the cross-product term $\sum_{i=1}^k \sum_{j=1}^{n_i} (\bar{Y}_{i.} - \bar{Y}_{..})(Y_{ij} - \bar{Y}_{i.})$ vanishes algebraically because $\sum_{j=1}^{n_i} (Y_{ij} - \bar{Y}_{i.}) = 0$. This yields the exact orthogonal partitioning known as the \textbf{Fundamental Sum of Squares Identity}:
\begin{align}
	\sum_{i=1}^k \sum_{j=1}^{n_i} (Y_{ij} - \bar{Y}_{..})^2 = \sum_{i=1}^k n_i (\bar{Y}_{i.} - \bar{Y}_{..})^2 + \sum_{i=1}^k \sum_{j=1}^{n_i} (Y_{ij} - \bar{Y}_{i.})^2.
\end{align}
Expressed symbolically:
\begin{align}
	\text{SCT} = \text{SCTR} + \text{SCE},
\end{align}
where:
\begin{itemize}
	\item $\text{SCT}$ (\textbf{Total Sum of Squares}): Measures the total variation across all $N$ data points relative to the grand mean across $\nu_T = N - 1$ degrees of freedom.
	\item $\text{SCTR}$ (\textbf{Treatment Sum of Squares}, also denoted $\text{SC}_{\text{Between}}$): Measures the systematic inter-group dispersion explained by differences across the $k$ treatment levels, with degrees of freedom $\nu_{\text{TR}} = k - 1$.
	\item $\text{SCE}$ (\textbf{Error Sum of Squares}, also denoted $\text{SC}_{\text{Within}}$ or $\text{Residual SS}$): Measures the unexplained intra-group variation within experimental units subjected to the exact same treatment, with degrees of freedom $\nu_E = \sum_{i=1}^k (n_i - 1) = N - k$.
\end{itemize}

Notice that the degrees of freedom partition exactly in the exact same additive manner as the sums of squares:
\begin{align}
	\nu_T = \nu_{\text{TR}} + \nu_E \implies (N - 1) = (k - 1) + (N - k).
\end{align}

\subsubsection{Expected Mean Squares and the Hypothesis Test}

A \textbf{Mean Square ($\text{CM}$)} is defined as a sum of squares divided by its associated degrees of freedom. Unlike sums of squares, mean squares are non-additive ($\text{CMT} \neq \text{CMTR} + \text{CME}$). We calculate:
\begin{align}
	\text{CMTR} &= \frac{\text{SCTR}}{k - 1}, \\
	\text{CME} &= \frac{\text{SCE}}{N - k}.
\end{align}
Taking mathematical expectations over the random model $Y_{ij} = \mu + \tau_i + \varepsilon_{ij}$ reveals the fundamental logic of ANOVA:
\begin{align}
	E[\text{CME}] &= \sigma^2, \\
	E[\text{CMTR}] &= \sigma^2 + \frac{\sum_{i=1}^k n_i \tau_i^2}{k - 1}.
\end{align}
Notice that $\text{CME}$ is an unconditionally unbiased estimator of the true experimental error variance $\sigma^2$. Conversely, $E[\text{CMTR}]$ equals $\sigma^2$ plus a non-negative quadratic term driven entirely by treatment effects ($\tau_i$).

We formulate the omnibus statistical hypotheses as:
\begin{align}
	H_0&: \tau_1 = \tau_2 = \dots = \tau_k = 0 \quad \Longleftrightarrow \quad \mu_1 = \mu_2 = \dots = \mu_k \\
	H_a&: \tau_i \neq 0 \text{ for at least one } i \quad \Longleftrightarrow \quad \mu_i \neq \mu_j \text{ for at least one pair } (i,j).
\end{align}
When the null hypothesis ($H_0$) is true, all treatment effects vanish ($\tau_i = 0$), implying that $E[\text{CMTR}] = \sigma^2$. Under $H_0$, both $\text{CMTR}$ and $\text{CME}$ represent independent estimates of the exact same underlying error variance $\sigma^2$. Therefore, by Cochran's Theorem, their ratio follows a central Snedecor $F$-distribution:
\begin{align}
	F_{\text{calc}} = \frac{\text{CMTR}}{\text{CME}} \sim F_{k-1, \, N-k}.
\end{align}
If the alternative hypothesis ($H_a$) is true, the quadratic term $\sum n_i \tau_i^2 / (k-1)$ is strictly positive, causing $E[\text{CMTR}] > E[\text{CME}]$. Consequently, the test statistic $F_{\text{calc}}$ will take systematically larger values, following a non-central $F$-distribution with non-centrality parameter $\lambda = \frac{1}{\sigma^2} \sum_{i=1}^k n_i \tau_i^2$.

\textbf{Decision Rule:} At a chosen significance level $\alpha$, we reject the null hypothesis $H_0$ if and only if the calculated statistic exceeds the upper critical quantile from statistical tables:
\begin{align}
	F_{\text{calc}} > F_{\alpha, \, k-1, \, N-k} \quad \text{or equivalently, if } p\text{-value} < \alpha.
\end{align}
All computational results are systematically compiled in the classic \textbf{ANOVA Table} shown in Table \ref{tab:anova_unifactorial}.

\begin{table}[h!]
	\centering
	\begin{tabular}{lcccc}
		\hline
		\textbf{Source of Variation} & \textbf{Sum of Squares (SC)} & \textbf{Degrees of Freedom ($\nu$)} & \textbf{Mean Square (CM)} & \textbf{$F$ Statistic} \\ \hline
		Treatments (Between) & $\text{SCTR} = \sum n_i(\bar{Y}_{i.} - \bar{Y}_{..})^2$ & $k - 1$ & $\text{CMTR} = \frac{\text{SCTR}}{k-1}$ & $F = \frac{\text{CMTR}}{\text{CME}}$ \\
		Error (Within) & $\text{SCE} = \sum \sum (Y_{ij} - \bar{Y}_{i.})^2$ & $N - k$ & $\text{CME} = \frac{\text{SCE}}{N-k}$ & \\ \hline
		\textbf{Total} & $\text{SCT} = \sum \sum (Y_{ij} - \bar{Y}_{..})^2$ & $N - 1$ & & \\ \hline
	\end{tabular}
	\caption{General Analysis of Variance Table for a One-Way Completely Randomized Design.}
	\label{tab:anova_unifactorial}
\end{table}

\subsection{Multiple Comparisons and Post-Hoc Tests}

When an omnibus ANOVA yields a statistically significant $F_{\text{calc}}$ ($p < \alpha$), we conclude that at least one treatment mean differs from the others ($\mu_i \neq \mu_j$). However, the omnibus test provides no specific information regarding \emph{which} pairs or combinations of groups are responsible for the rejection. To pinpoint the exact location of significant differences without inflating the overall family-wise error rate ($\text{FWER}$), specialized \textbf{Post-Hoc Multiple Comparison Procedures} must be applied.

\subsubsection{Tukey's Honestly Significant Difference (HSD)}
\textbf{Tukey's HSD test} is the most widely recommended and statistically optimal method for conducting all possible $\binom{k}{2} = \frac{k(k-1)}{2}$ pairwise comparisons across balanced treatment groups ($n_1 = n_2 = \dots = n_k = n$). It rigorously controls the exact family-wise error rate at $\alpha$ across the entire set of pairwise contrasts.

Tukey's procedure is derived from the \textbf{Studentized Range Distribution ($q$)}, which describes the sampling distribution of the difference between the sample maximum and sample minimum ($\bar{Y}_{\text{max}} - \bar{Y}_{\text{min}}$) divided by the pooled standard error.

The \textbf{Honestly Significant Difference ($\text{HSD}$)} critical margin for comparing any two treatment sample means $\bar{Y}_{i.}$ and $\bar{Y}_{j.}$ is computed as:
\begin{align}
	\text{HSD} = q_{\alpha, \, k, \, \nu_E} \sqrt{\frac{\text{CME}}{n}},
\end{align}
where $q_{\alpha, \, k, \, \nu_E}$ is the upper critical quantile from the Studentized range table with $k$ treatments and $\nu_E = N - k$ error degrees of freedom, and $\text{CME}$ is the error mean square from the ANOVA table. If group sample sizes are unbalanced ($n_i \neq n_j$), the \textbf{Tukey-Kramer adjustment} replaces $n$ with the harmonic mean of the two group sample sizes:
\begin{align}
	\text{HSD}_{ij} = \frac{q_{\alpha, \, k, \, \nu_E}}{\sqrt{2}} \sqrt{\text{CME} \left( \frac{1}{n_i} + \frac{1}{n_j} \right)}.
\end{align}

\textbf{Decision Rule for Tukey's HSD:} For any pair of treatments $(i, j)$, we declare that population means $\mu_i$ and $\mu_j$ are statistically different at overall family level $\alpha$ if and only if the absolute difference between their sample means exceeds the calculated $\text{HSD}$ threshold:
\begin{align}
	|\bar{Y}_{i.} - \bar{Y}_{j.}| > \text{HSD} \implies \text{Reject } H_0: \mu_i = \mu_j.
\end{align}

\subsubsection{Fisher's Least Significant Difference (LSD)}
\textbf{Fisher's LSD test} is essentially a direct extension of the standard two-sample Student's $t$-test across all pairs of means, utilizing the pooled error mean square ($\text{CME}$) from the ANOVA across $\nu_E = N - k$ degrees of freedom. The critical difference threshold between any two sample means is defined as:
\begin{align}
	\text{LSD} = t_{\alpha/2, \, \nu_E} \sqrt{\text{CME} \left( \frac{1}{n_i} + \frac{1}{n_j} \right)}.
\end{align}
\textbf{Critical Limitation:} Fisher's LSD does not incorporate a family-wise correction factor across multiple tests. Consequently, it suffers from severe Type I error inflation when $k > 3$. To mitigate false positives, Fisher proposed \textbf{Protected LSD}, which mandates that LSD comparisons may only be computed if the omnibus $F$-test is statistically significant at $\alpha$. While Protected LSD successfully controls $\text{FWER}$ at $\alpha$ when $k=3$, it fails completely when $k \geq 4$. Therefore, Fisher's LSD is primarily utilized for exploratory research or when statistical sensitivity (high power) is prioritized over Type I error protection.

\subsubsection{Bonferroni Correction}
The \textbf{Bonferroni correction} is a universal, simple, and conservative post-hoc method based on Boole's inequality for intersection event probabilities. If the researcher plans to perform $m = \binom{k}{2}$ pairwise comparisons and wishes to maintain the overall family-wise confidence level at or above $(1-\alpha)100\%$, the individual significance level for each specific pairwise test is adjusted by dividing the nominal level by the total number of comparisons:
\begin{align}
	\alpha^* = \frac{\alpha}{m} = \frac{\alpha}{\binom{k}{2}}.
\end{align}
The critical threshold for declaring two treatment means significantly different using the Bonferroni procedure becomes:
\begin{align}
	\text{CD}_{\text{Bonf}} = t_{\alpha^*/2, \, \nu_E} \sqrt{\text{CME} \left( \frac{1}{n_i} + \frac{1}{n_j} \right)}.
\end{align}
While Bonferroni absolutely prevents false-positive rejections ($\text{FWER} \leq \alpha$), it becomes excessively conservative when the number of treatments $k$ is large ($k \geq 6 \implies m \geq 15$), drastically increasing Type II errors ($\beta$) and sacrificing statistical power.

\subsection{Randomized Complete Block Design (RCBD)}

When experimental units possess inherent, known heterogeneity (such as different batches of raw material, distinct age cohorts, or different physical testing machines), applying a Completely Randomized Design (CRD) forces all inter-unit variability directly into the residual error term ($\text{SCE}$). This inflates $\text{CME}$ and destroys the sensitivity of the $F$-test.

The \textbf{Randomized Complete Block Design (RCBD)} overcomes this challenge by grouping experimental units into $b$ homogeneous blocks of size $k$ (where $k$ is the total number of treatments). Within each block, treatments are assigned completely at random, ensuring that every treatment appears exactly once per block ($N = kb$).

\subsubsection{Additive Structural Model for RCBD}
The observation $Y_{ij}$ for the $i$-th treatment ($i=1,\dots,k$) within the $j$-th block ($j=1,\dots,b$) is modeled by:
\begin{align}
	Y_{ij} = \mu + \tau_i + \beta_j + \varepsilon_{ij},
\end{align}
where:
\begin{itemize}
	\item $\mu$ is the overall population grand mean.
	\item $\tau_i$ is the fixed effect of the $i$-th treatment ($\sum_{i=1}^k \tau_i = 0$).
	\item $\beta_j$ is the effect of the $j$-th block ($\sum_{j=1}^b \beta_j = 0$), capturing systematic shifts across blocks.
	\item $\varepsilon_{ij} \sim N(0, \sigma^2)$ are independent random residual errors, assuming no interaction between treatments and blocks (additivity of effects).
\end{itemize}

\subsubsection{Partitioning of Variance in RCBD}
In an RCBD, the Total Sum of Squares ($\text{SCT}$, with $\nu_T = kb - 1$) is orthogonally partitioned into three mutually independent components:
\begin{align}
	\text{SCT} = \text{SCTR} + \text{SCB} + \text{SCE},
\end{align}
where:
\begin{itemize}
	\item $\text{SCTR} = b \sum_{i=1}^k (\bar{Y}_{i.} - \bar{Y}_{..})^2$, with $\nu_{\text{TR}} = k - 1$.
	\item $\text{SCB} = k \sum_{j=1}^b (\bar{Y}_{.j} - \bar{Y}_{..})^2$, with $\nu_{\text{B}} = b - 1$, measuring structural variability isolated between blocks ($\bar{Y}_{.j} = \frac{1}{k} \sum_{i=1}^k Y_{ij}$).
	\item $\text{SCE} = \sum_{i=1}^k \sum_{j=1}^b (Y_{ij} - \bar{Y}_{i.} - \bar{Y}_{.j} + \bar{Y}_{..})^2$, with residual degrees of freedom $\nu_E = (k - 1)(b - 1)$.
\end{itemize}

By comparing the expected mean squares:
\begin{align}
	E[\text{CME}] &= \sigma^2, \\
	E[\text{CMTR}] &= \sigma^2 + \frac{b \sum_{i=1}^k \tau_i^2}{k - 1}, \\
	E[\text{CMB}] &= \sigma^2 + \frac{k \sum_{j=1}^b \beta_j^2}{b - 1},
\end{align}
we construct two independent $F$-statistics: one to evaluate treatment differences ($F_{\text{TR}} = \frac{\text{CMTR}}{\text{CME}} \sim F_{k-1, \, (k-1)(b-1)}$) and another to assess the statistical effectiveness of the blocking factor ($F_{\text{B}} = \frac{\text{CMB}}{\text{CME}} \sim F_{b-1, \, (k-1)(b-1)}$), as summarized in Table \ref{tab:anova_dbca}.

\begin{table}[h!]
	\centering
	\begin{tabular}{lcccc}
		\hline
		\textbf{Source of Variation} & \textbf{Sum of Squares (SC)} & \textbf{d.f. ($\nu$)} & \textbf{Mean Square (CM)} & \textbf{$F$ Statistic} \\ \hline
		Treatments & $\text{SCTR} = b \sum (\bar{Y}_{i.} - \bar{Y}_{..})^2$ & $k - 1$ & $\text{CMTR} = \frac{\text{SCTR}}{k-1}$ & $F_{\text{TR}} = \frac{\text{CMTR}}{\text{CME}}$ \\
		Blocks & $\text{SCB} = k \sum (\bar{Y}_{.j} - \bar{Y}_{..})^2$ & $b - 1$ & $\text{CMB} = \frac{\text{SCB}}{b-1}$ & $F_{\text{B}} = \frac{\text{CMB}}{\text{CME}}$ \\
		Residual Error & $\text{SCE} = \text{SCT} - \text{SCTR} - \text{SCB}$ & $(k-1)(b-1)$ & $\text{CME} = \frac{\text{SCE}}{(k-1)(b-1)}$ & \\ \hline
		\textbf{Total} & $\text{SCT} = \sum \sum (Y_{ij} - \bar{Y}_{..})^2$ & $kb - 1$ & & \\ \hline
	\end{tabular}
	\caption{ANOVA Table for a Randomized Complete Block Design (RCBD).}
	\label{tab:anova_dbca}
\end{table}
